\documentclass[10pt]{beamer}

% --- Design-Einstellungen ---
\usetheme{metropolis} % Modernes, cleanes Design
\usecolortheme{owl}     % Dark-Mode (passend zum Hackintosh-Thema)
\usepackage[utf8]{inputenc}
\usepackage[ngerman]{babel}
\usepackage{graphicx}
\usepackage{booktabs}
\usepackage{listings}

% --- Metadaten ---
\title{Hackintosh Companion App}
\subtitle{Mobile Begleit-Applikation zur Systemkonfiguration}
\author{Kilian Ketelhohn}
\institute{Hochschule Wismar \\ Fachbereich Informatik}
\date{Verteidigung: 11. April 2026}

\begin{document}

% Titelfolie
\begin{frame}
    \titlepage
\end{frame}

% Inhaltsverzeichnis
\begin{frame}{Agenda}
    \tableofcontents
\end{frame}

\section{Problemstellung \& Motivation}
\begin{frame}{Motivation}
    \begin{itemize}
        \item \textbf{Die Herausforderung:} Hackintosh-Konfiguration ist komplex (ACPI, Kexts, OpenCore).
        \item \textbf{Das Problem:} Dokumentation ist oft nur online verfügbar, während das Zielgerät offline ist.
        \item \textbf{Die Lösung:} Eine mobile, offline-fähige Wissensdatenbank mit interaktiven Features.
    \end{itemize}
\end{frame}

\section{Technologie-Stack}
\begin{frame}{Entscheidungen & Architektur}
    \begin{columns}
        \begin{column}{0.5\textwidth}
            \textbf{Frontend/Engine:}
            \begin{itemize}
                \item Flutter (Cross-Platform)
                \item Dart (AOT-kompiliert)
                \item Skia-Renderer (Legacy-Support)
            \end{itemize}
        \end{column}
        \begin{column}{0.5\textwidth}
            \textbf{Datenhaltung:}
            \begin{itemize}
                \item SQLite (Lokal)
                \item REST-API (Cloud-Sync)
                \item JSON-Mapping
            \end{itemize}
        \end{column}
    \end{columns}
\end{frame}

\section{Kern-Features}
\begin{frame}{Technische Highlights}
    \begin{block}{1. Boot-Simulation}
        Visualisierung des Startvorgangs (BIOS/macOS) zur UX-Steigerung.
    \end{block}
    \begin{block}{2. 3D-Hardware-Visualisierung}
        Interaktives Laptop-Modell zur Komponenten-Identifikation.
    \end{block}
    \begin{block}{3. QR-Code-Integration}
        Fehlerfreie Übertragung von Konfigurations-Strings via Kamera.
    \end{block}
\end{frame}

\section{Implementierung}
\begin{frame}[fragile]{Offline-First Logik (Auszug)}
\begin{lstlisting}[language=SQL, basicstyle=\tiny]
CREATE TABLE devices (
  id INTEGER PRIMARY KEY AUTOINCREMENT,
  name TEXT,
  cpu TEXT,
  compatible INTEGER
);
\end{lstlisting}
\begin{itemize}
    \item \textbf{Single-Source-of-Truth:} Lokale DB steht immer an erster Stelle.
    \item \textbf{Asynchrone Kommunikation:} API-Sync im Hintergrund.
\end{itemize}
\end{frame}

\section{Evaluation}
\begin{frame}{Performance auf Samsung A3 (2016)}
    \begin{table}
        \centering
        \begin{tabular}{lr}
            \toprule
            Metrik & Messwert \\
            \midrule
            App-Startzeit & $\approx$ 2,5 s \\
            DB-Abfrage (50 Einträge) & < 15 ms \\
            3D-Renderer (Skia) & 55--60 FPS \\
            QR-Erkennung & < 1,5 s \\
            \bottomrule
        \end{tabular}
    \end{table}
    \textbf{Ergebnis:} Erfolgreicher Nachweis der Praxistauglichkeit auf Legacy-Hardware.
\end{frame}

\section{Fazit}
\begin{frame}{Zusammenfassung \& Ausblick}
    \begin{itemize}
        \item \textbf{Erreicht:} Voll funktionsfähiges MVP (Minimum Viable Product).
        \item \textbf{Lerneffekt:} Tiefer Einblick in Grafik-Pipelines und mobile Daten-Synchronisation.
        \item \textbf{Ausblick:} GitHub-API Anbindung für automatisierte Kext-Updates.
    \end{itemize}
\end{frame}

% Abschlussfolie
\begin{frame}
    \centering
    \Huge Vielen Dank für Ihre Aufmerksamkeit! \\
    \vspace{1cm}
    \Large Fragen?
\end{frame}

\end{document}